% GENERATED FILE. Edit canonical lessons, then run npm run generate:books.
% canonical-source-hash: fnv1a64:850b7fce319f53f8
% canonical-lessons: FA-C299-agar-vaqt-dashtam, FA-C299-dust-dashtam, FA-R299-would-and-if-only, FA-W299-shesh-khat

\chapter{Would, And Six Lines About A Trip}
\label{ch:fa-would-and-six-lines}
\hlchaptermodality{299}

\begin{chapteropening}
\textbf{By the end of this chapter:} \emph{I can say what I would do if things were different, say what I would like, and write six connected lines about a trip.}

Complete FA-W299-shesh-khat: six connected lines to one reader about a trip, with a scene, events in order, a reason and a condition.
\end{chapteropening}

\section[\texorpdfstring{\fa{اگر} \fa{وقت} \fa{داشتم،} … (agar vaqt …)}{اگر وقت داشتم، … (agar vaqt …)}]{\fa{اگر} \fa{وقت} \fa{داشتم،} \fa{میرفتم} (agar vaqt dâshtam, miraftam) — if I had time, I would go}

\label{lesson:FA-C299-agar-vaqt-dashtam}



\begin{quote}
Say \textquotedblleft{}if the weather is good, we'll go\textquotedblright{}. (\emph{Agar havâ khub bâshad, miravim.}) Then \textquotedblleft{}it depends\textquotedblright{}. (\emph{Bastegi dârad.})

That condition may come true. The last turn in this run is a condition that is not true at all, and what would follow if it were.
\end{quote}

\begin{grammarlens}[title={Grammar Lens: the plain past after agar, the mi- past after it}]
\begin{quote}
\textbf{\fa{اگر} \fa{وقت} \fa{داشتم،} \fa{به} \fa{سفر} \fa{میرفتم}.} — \emph{agar vaqt dâshtam, be safar miraftam} — If I had time, I would go on a trip.
\end{quote}

I have no time, and I am not going. The \emph{agar} half says so by moving into the plain past, \emph{dâshtam}, the way English moves to \textquotedblleft{}had\textquotedblright{}. The result goes into the \emph{mi-} past, \emph{miraftam}.

Set it beside the real one: \emph{agar havâ khub bâshad, miravim}, it may be fine, and then we'll go; \emph{agar havâ khub bud, miraftim}, it is not fine, so we are not going.
\end{grammarlens}

\begin{grammarlens}[title={Grammar Lens: the mi- past as \textquotedblleft{}would\textquotedblright{}}]
The shape from chapter two hundred and ninety, \emph{miraftam}, \textquotedblleft{}I was going\textquotedblright{}, does a third job here: \textquotedblleft{}I would go\textquotedblright{}.

\begin{quote}
\textbf{\fa{اگر} \fa{پول} \fa{داشتم،} \fa{یک} \fa{ماشین} \fa{میخریدم}.} — \emph{agar pul dâshtam, yek mâshin mikharidam} — If I had the money, I'd buy a car.
\end{quote}

\begin{quote}
\textbf{\fa{اگر} \fa{جای} \fa{تو} \fa{بودم،} \fa{زودتر} \fa{میخوابیدم}.} — \emph{agar jâ-ye to budam, zudtar mikhâbidam} — If I were you, I'd go to bed earlier.
\end{quote}

\emph{Jâ-ye to} is \textquotedblleft{}your place\textquotedblright{}, so the advice is \textquotedblleft{}if I were in your place\textquotedblright{}. \emph{Budam} and \emph{dâshtam} keep their plain past here too, with no \emph{mi-}.

The same pair can also look back, \textquotedblleft{}if I had had time, I would have gone\textquotedblright{}: \emph{agar vaqt dâshtam, miraftam} covers both. The situation tells you which.
\end{grammarlens}

\subsection*{Your turn}
Say these aloud:

\begin{itemize}
\raggedright
  \item \emph{agar vaqt dâshtam, be safar miraftam}
  \item \emph{agar pul dâshtam, yek mâshin mikharidam}
  \item advice — \emph{agar jâ-ye to budam, zudtar mikhâbidam}
  \item real, then unreal — \emph{agar havâ khub bâshad, miravim}; \emph{agar havâ khub bud, miraftim}
\end{itemize}

\begin{tcolorbox}[breakable,colback=teal!4,colframe=teal!35!black,title={Before you move on}]
Which past goes after \emph{agar} when the condition is not true? (The \textbf{plain past}: \emph{dâshtam}.) Which says \textquotedblleft{}would\textquotedblright{}? (The \emph{\textbf{mi-}} past: \emph{miraftam}.) \textquotedblleft{}If I were you\textquotedblright{}? (\emph{\textbf{Agar jâ-ye to budam.}})
\end{tcolorbox}
\section[\texorpdfstring{\fa{دوست} \fa{داشتم} (dust dâshtam)}{دوست داشتم (dust dâshtam)}]{\fa{دوست} \fa{داشتم} (dust dâshtam) — I would like, and if only}

\label{lesson:FA-C299-dust-dashtam}



\begin{quote}
Say \textquotedblleft{}I'd like to go\textquotedblright{}. (\emph{Dust dâram beravam.}) Then \textquotedblleft{}if I had time, I would go\textquotedblright{}. (\emph{Agar vaqt dâshtam, miraftam.})

Put those two together and you have the politest way to turn something down.
\end{quote}

\begin{grammarlens}[title={Grammar Lens: dust dâshtam, for a wish held back}]
\emph{Dust dâram beravam} is open: I'd like to go, and I may. Move \emph{dâshtan} into the past and the wish steps back from what is happening:

\begin{quote}
\textbf{\fa{دوست} \fa{داشتم} \fa{بیایم،} \fa{ولی} \fa{وقت} \fa{ندارم}.} — \emph{dust dâshtam biyâyam, vali vaqt nadâram} — I'd love to come, but I don't have time.
\end{quote}

The wished-for verb stays in the subjunctive, \emph{biyâyam}, as it did after \emph{dust dâram}. Only \emph{dâshtan} moves. Said about the past, it is \textquotedblleft{}I would have liked\textquotedblright{}: \emph{dust dâshtam ân film râ bebinam}, \textquotedblleft{}I'd have liked to see that film\textquotedblright{}. \emph{Bebinam} is \emph{didan} in the subjunctive, on its stem \emph{bin-}.
\end{grammarlens}

\begin{grammarlens}[title={Grammar Lens: kâsh, \textquotedblleft{}if only\textquotedblright{}}]
\textbf{\fa{کاش}}, \emph{kâsh}, is a wish with no result attached. With the plain past it wishes against what is true now:

\begin{quote}
\textbf{\fa{کاش} \fa{وقت} \fa{داشتم}.} — \emph{kâsh vaqt dâshtam} — If only I had time.
\end{quote}

With the subjunctive it wishes for something still open:

\begin{quote}
\textbf{\fa{کاش} \fa{فردا} \fa{هوا} \fa{خوب} \fa{باشد}.} — \emph{kâsh fardâ havâ khub bâshad} — I do hope the weather is good tomorrow.
\end{quote}

That second one sits close to \emph{omidvâram ke}, with more longing in it.
\end{grammarlens}

\subsection*{Your turn}
Say these aloud:

\begin{itemize}
\raggedright
  \item turn an invitation down kindly — \emph{dust dâshtam biyâyam, vali vaqt nadâram}
  \item \textquotedblleft{}If only I had time\textquotedblright{} — \emph{kâsh vaqt dâshtam}
  \item an open wish — \emph{kâsh fardâ havâ khub bâshad}
  \item open, then held back — \emph{dust dâram beravam}, \emph{dust dâshtam beravam}
\end{itemize}

\begin{tcolorbox}[breakable,colback=teal!4,colframe=teal!35!black,title={Before you move on}]
Which is open, \emph{dust dâram} or \emph{dust dâshtam}? (\emph{\textbf{Dust dâram.}}) Which verb stays in the subjunctive after either? (The wished-for one: \emph{biyâyam}.) \textquotedblleft{}If only I had time\textquotedblright{}? (\emph{\textbf{Kâsh vaqt dâshtam.}})
\end{tcolorbox}
\section[(review)]{Review — what would be, and what may be}

\label{lesson:FA-R299-would-and-if-only}



\begin{quote}
Say \textquotedblleft{}if I had the money, I'd buy a car\textquotedblright{}. (\emph{Agar pul dâshtam, yek mâshin mikharidam.})
\end{quote}

\begin{grammarlens}[title={What you've built}]
Two chapters, and two kinds of \textquotedblleft{}if\textquotedblright{}. A condition that may come true takes the subjunctive: \emph{agar bârân biyâyad, dar khâne mimânim}. One that is not true takes the plain past, and its result the \emph{mi-} past: \emph{agar havâ khub bud, be pârk miraftim}, \textquotedblleft{}if the weather were good, we'd go to the park\textquotedblright{}.

The \emph{mi-} past has now done three jobs. In chapter two hundred and ninety it painted a scene, \emph{mikhândam}, and told a habit, \emph{har ruz miraftam}. Here it says what would happen, \emph{miraftim}.

\emph{Dust dâshtam} holds a wish back, and \emph{kâsh} sighs for one:

\begin{quote}
\emph{dust dâshtam be safar beravam, vali vaqt nadâram. kâsh vaqt dâshtam!}
\end{quote}
\end{grammarlens}

\subsection*{Your turn}
\begin{itemize}
\raggedright
  \item \emph{Say it:} may be, then is not — \emph{agar bârân biyâyad, dar khâne mimânim}; \emph{agar havâ khub bud, be pârk miraftim}
  \item \emph{Say it:} \emph{dust dâshtam be safar beravam, vali vaqt nadâram. kâsh vaqt dâshtam!}
  \item \emph{Answer:} \emph{fardâ be pârk miravi?} — with \emph{bastegi dârad} and a real condition
  \item \emph{Say it:} advice with \emph{agar jâ-ye to budam} and a cardigan — \emph{... zhâkat mipushidam}
\end{itemize}

\begin{tcolorbox}[breakable,colback=teal!4,colframe=teal!35!black,title={Before you move on}]
\textquotedblleft{}If it rains\textquotedblright{} and \textquotedblleft{}if the weather were good\textquotedblright{}: which verb follows \emph{agar} in each? (\emph{\textbf{Biyâyad}}, the subjunctive; \emph{\textbf{bud}}, the plain past.) \textquotedblleft{}If only\textquotedblright{}? (\emph{\textbf{Kâsh.}}) Which letter opens \textbf{\fa{ژاکت}}? (\textbf{\fa{ژ}}, \emph{zhe}.)
\end{tcolorbox}
\section[(six lines)]{Six lines about a trip}

\label{lesson:FA-W299-shesh-khat}



\begin{quote}
Chapter two hundred and eighty-nine asked for four lines that held together. This one asks for six, and they tell a story.

Pick one reader: a friend, by name.
\end{quote}

\subsection*{What to know first}
A told story is not a list of events. It has a \textbf{scene}, in \emph{bud} and the \emph{mi-} past. It has \textbf{events} in the plain past, put in order with \emph{avval}, \emph{ba'd}, \emph{vaqti} and \emph{dar âkhar}. It has at least one \textbf{reason}, with \emph{chon} or \emph{barâ-ye hamin}, and a \textbf{feeling}, with \emph{budam} or \emph{shodam}.

A good one ends by looking forward or wishing: \emph{omidvâram ke...}, \emph{agar ... bâshad, ...}, or \emph{agar vaqt dâshtam, ... miraftam}.

The B1 writing paper asks for a connected text of 130 to 170 words. Six lines is the first step towards it, not the whole of it.

\subsection*{Writing — connected composition}
On paper, right to left, write \textbf{six lines} to your reader about a trip, real or invented, one line for each of these:

\begin{enumerate}
\raggedright
  \item the scene: when it was, with \emph{gozashte} or \emph{pish}, and the weather;
  \item where you went, who was with you, with \emph{bâ}, and for how long;
  \item the first thing that happened, with \emph{avval};
  \item the next thing, with \emph{ba'd} or \emph{vaqti}, and a reason, with \emph{chon}, \emph{be khâter-e} or \emph{barâ-ye hamin};
  \item how it ended and how you felt, with \emph{dar âkhar} and \emph{budam} or \emph{shodam};
  \item one line forward: \emph{omidvâram ke}, a real \emph{agar}, or \emph{agar vaqt dâshtam}.
\end{enumerate}

Use only words you already write. The whole exercise is choosing the right past for each line, and joining them.

\begin{tcolorbox}[breakable,colback=teal!4,colframe=teal!35!black,title={Before you move on}]
Now read your lines as you read the winter trip in chapter two hundred and ninety-one. Mark every verb \textbf{S} for scene or \textbf{E} for event. Line one should be mostly \textbf{S}; lines three and four mostly \textbf{E}. A line about the weather marked \textbf{E} is the commonest slip: change \emph{sard shod} back to \emph{sard bud}.

Then three spelling checks. Every \emph{mi-} is joined to its verb, \textbf{\fa{میرفتم}}. Every \textbf{\fa{و}} is followed by a break smaller than the space between words. And each word keeps its own letter: \textbf{\fa{غذا}} with \textbf{\fa{غ}}, \textbf{\fa{قطار}} with \textbf{\fa{ق}}, though the two sound alike.

Last, read line six alone. If it has \emph{agar} and something that may happen, is the verb after \emph{agar} in the subjunctive? If it is a wish against the facts, is it in the plain past, with the \emph{mi-} past after it?
\end{tcolorbox}
